\documentclass[10pt,a4paper]{amsart}
\usepackage{amsmath,amssymb,amsthm,graphicx}
\usepackage{verbatim}
\newenvironment{reference}{\comment}{\endcomment}
%\newenvironment{reference}{}{}
\newenvironment{slogan}{\comment}{\endcomment}
\newenvironment{history}{\comment}{\endcomment}

% For commutative diagrams we use Xy-pic
\usepackage[all]{xy}

% We use 2cell for 2-commutative diagrams.
\xyoption{2cell}
% Theorem environments.
%
\theoremstyle{plain}
\newtheorem{theorem}[subsection]{Theorem}
\newtheorem{proposition}[subsection]{Proposition}
\newtheorem{lemma}[subsection]{Lemma}

\theoremstyle{definition}
\newtheorem{definition}[subsection]{Definition}
\newtheorem{example}[subsection]{Example}
\newtheorem{exercise}[subsection]{Exercise}
\newtheorem{situation}[subsection]{Situation}

\theoremstyle{remark}
\newtheorem{remark}[subsection]{Remark}
\newtheorem{remarks}[subsection]{Remarks}

\numberwithin{equation}{subsection}

% Macros
%
\def\lim{\mathop{\mathrm{lim}}\nolimits}
\def\colim{\mathop{\mathrm{colim}}\nolimits}
\def\Spec{\mathop{\mathrm{Spec}}}
\def\Hom{\mathop{\mathrm{Hom}}\nolimits}
\def\Ext{\mathop{\mathrm{Ext}}\nolimits}
\def\SheafHom{\mathop{\mathcal{H}\!\mathit{om}}\nolimits}
\def\SheafExt{\mathop{\mathcal{E}\!\mathit{xt}}\nolimits}
\def\Sch{\mathit{Sch}}
\def\Mor{\mathop{\mathrm{Mor}}\nolimits}
\def\Ob{\mathop{\mathrm{Ob}}\nolimits}
\def\Sh{\mathop{\mathit{Sh}}\nolimits}
\def\NL{\mathop{N\!L}\nolimits}
\def\CH{\mathop{\mathrm{CH}}\nolimits}
\def\proetale{{pro\text{-}\acute{e}tale}}
\def\etale{{\acute{e}tale}}
\def\QCoh{\mathit{QCoh}}
\def\Ker{\mathop{\mathrm{Ker}}}
\def\Im{\mathop{\mathrm{Im}}}
\def\Coker{\mathop{\mathrm{Coker}}}
\def\Coim{\mathop{\mathrm{Coim}}}

% Boxtimes
%
\DeclareMathSymbol{\boxtimes}{\mathbin}{AMSa}{"02}

%
% Macros for moduli stacks/spaces
%
\def\QCohstack{\mathcal{QC}\!\mathit{oh}}
\def\Cohstack{\mathcal{C}\!\mathit{oh}}
\def\Spacesstack{\mathcal{S}\!\mathit{paces}}
\def\Quotfunctor{\mathrm{Quot}}
\def\Hilbfunctor{\mathrm{Hilb}}
\def\Curvesstack{\mathcal{C}\!\mathit{urves}}
\def\Polarizedstack{\mathcal{P}\!\mathit{olarized}}
\def\Complexesstack{\mathcal{C}\!\mathit{omplexes}}
% \Pic is the operator that assigns to X its picard group, usage \Pic(X)
% \Picardstack_{X/B} denotes the Picard stack of X over B
% \Picardfunctor_{X/B} denotes the Picard functor of X over B
\def\Pic{\mathop{\mathrm{Pic}}\nolimits}
\def\Picardstack{\mathcal{P}\!\mathit{ic}}
\def\Picardfunctor{\mathrm{Pic}}
\def\Deformationcategory{\mathcal{D}\!\mathit{ef}}
\usepackage{cleveref}
\crefname{originallemma}{original lemma}{original lemmas}
\Crefname{originallemma}{Original lemma}{Original lemmas}
\makeatletter
\newlabel{lemma-lisse-etale}{{0788}{source}{}{sourcepointer-1}{}}
\newlabel{lemma-lisse-etale@cref}{{[originallemma][0][]0788}{[1][0][]source}{}{}{}}
\makeatother
\title{Parasitic modules detected on the lisse-etale and flat-fppf sites}
\author{The Stacks Project Authors}
\date{stacks-project:stacks-cohomology.tex:local-snapshot-20261012; GNU FDL 1.2 or later}
\begin{document}
\maketitle
\noindent\textit{Editorial note.} All delivered mathematical text is the author's own, copied verbatim from the pinned Stacks Project chapter file stacks-cohomology.tex; no statement or proof step was rewritten, reordered or removed.\par\medskip
\noindent\textit{Editorial note.} The counted span cites the lemmas recorded in assembly.external\_references with canonical tags as navigation pointers.\par\medskip
\noindent\textit{Allowed standard prerequisites (editorial).}\begin{itemize}
\item Faithfulness of smooth surjective pullback on qc modules. Standard fact used in this chapter (also in adjoint-kernel-parasitic). q\_\{small\}\textasciicircum{}* faithful.
\item Flat base change property. Cohomology of Algebraic Stacks, Definition stacks-cohomology-definition-flat-base-change. p\textasciicircum{}*(F|\_U) = q\textasciicircum{}*(F|\_V).
\end{itemize}\medskip
\ifdefined\hypertarget\hypertarget{sourcepointer-1}{}\fi
\noindent\textit{Source reference.} Original lemma 0788: sources/stacks-chow/tags.txt, source lines 9228--9228. Cross-chapter standard prerequisite cited by the author inside the counted span; it is used as a pointer and its own proof is not part of this item.\par\medskip
\begin{lemma}
\label{lemma-parasitic-in-terms-flat-fppf}
Let $\mathcal{X}$ be an algebraic stack.
\begin{enumerate}
\item Let $\mathcal{F}$ be an $\mathcal{O}_\mathcal{X}$-module
with the flat base change property on $\mathcal{X}_\etale$.
The following are equivalent
\begin{enumerate}
\item $\mathcal{F}$ is parasitic, and
\item $g^*\mathcal{F} = 0$ where
$g : \Sh(\mathcal{X}_{lisse,\etale}) \to
\Sh(\mathcal{X}_\etale)$ is as in Lemma \ref{lemma-lisse-etale}.
\end{enumerate}
\item Let $\mathcal{F}$ be an $\mathcal{O}_\mathcal{X}$-module on
$\mathcal{X}_{fppf}$. The following are equivalent
\begin{enumerate}
\item $\mathcal{F}$ is parasitic, and
\item $g^*\mathcal{F} = 0$ where
$g :  \Sh(\mathcal{X}_{flat,fppf}) \to \Sh(\mathcal{X}_{fppf})$
is as in Lemma \ref{lemma-lisse-etale}.
\end{enumerate}
\end{enumerate}
\end{lemma}

\begin{proof}
Part (2) is immediate from the definitions (this is one of the advantages
of the flat-fppf site over the lisse-\'etale site). The implication
(1)(a) $\Rightarrow$ (1)(b) is immediate as well. To see (1)(b)
$\Rightarrow$ (1)(a) let $U$ be a scheme and let $x : U \to \mathcal{X}$
be a surjective smooth morphism. Then $x$ is an object of the
lisse-\'etale site of $\mathcal{X}$. Hence we see that (1)(b)
implies that $\mathcal{F}|_{U_\etale} = 0$. Let $V \to \mathcal{X}$
be an flat morphism where $V$ is a scheme. Set $W = U \times_\mathcal{X} V$
and consider the diagram
$$
\xymatrix{
W \ar[d]_p \ar[r]_q & V \ar[d] \\
U \ar[r] & \mathcal{X}
}
$$
Note that the projection $p : W \to U$ is flat and the projection
$q : W \to V$ is smooth and surjective. This implies that $q_{small}^*$
is a faithful functor on quasi-coherent modules. By assumption $\mathcal{F}$
has the flat base change property so that we obtain
$p_{small}^*\mathcal{F}|_{U_\etale} \cong
q_{small}^*\mathcal{F}|_{V_\etale}$. Thus if $\mathcal{F}$
is in the kernel of $g^*$, then $\mathcal{F}|_{V_\etale} = 0$
as desired.
\end{proof}

\end{document}
